Srovnávací tabulka \ref{tab:flexicurity} shrnuje klíčové indikátory trhu práce pro šest zemí střední a~severní Evropy za rok 2023 a~odhaluje strukturální paradox české ekonomiky: ČR dosahuje vysoké zaměstnanosti (\SI{81.7}{\percent}), nízkého \ac{AROPE} (\SI{12.0}{\percent}) a~relativně komprimovaného Giniho koeficientu (24,4) --- avšak za cenu nejvyššího počtu odpracovaných hodin v~souboru (\SI{39.6}{\hour\per\week}), daňového klínu přesahujícího Dánsko (\SI{38.1}{\percent} vs.~\SI{33.5}{\percent}), a~zároveň výrazně nízkého pokrytí kolektivními smlouvami (\SI{32}{\percent}).
Dánský model --- nejnižší pracovní doba, nejnižší \ac{AROPE} bez extrémního Giniho --- je dosažen při vysoké odborové hustotě (\SI{67}{\percent}) a~\ac{KS} pokrytí (\SI{80}{\percent}): demonstruje, že flexibilita trhu práce a~sociální ochrana nejsou protiklady, pokud jsou vyváženy silnými institucemi kolektivního vyjednávání.
Rakousco kombinuje nejvyšší \ac{KS} pokrytí (\SI{98}{\percent}, erga omnes) s~nízkým \ac{AROPE} (\SI{17.7}{\percent}) a~nadprůměrnou mzdou: ukazuje, že i~bez skandinávské odborové hustoty lze dosáhnout solidní mzdové ochrany, pokud je rozšiřování \ac{KS} právně garantováno.
ČR a~Slovensko --- oba s~nízkým pokrytím --- pracují více hodin, vydělávají méně v~\ac{PPS} a~mají vyšší daňový klín na pracovníka: tento trojúhelník naznačuje, že česká konkurenceschopnost je stále postavena na modelu levné práce, nikoliv na kvalitní institucionální infrastruktuře trhu práce.

Srovnávací tabulka \ref{tab:flexicurity} shrnuje klíčové indikátory trhu práce pro šest zemí střední a~severní Evropy za rok 2023 a~odhaluje strukturální paradox české ekonomiky: ČR dosahuje vysoké zaměstnanosti (\SI{81.7}{\percent}), nízkého \ac{AROPE} (\SI{12.0}{\percent}) a~relativně komprimovaného Giniho koeficientu (24,4) --- avšak za cenu nejvyššího počtu odpracovaných hodin v~souboru (\SI{39.6}{\hour\per\week}), daňového klínu přesahujícího Dánsko (\SI{38.1}{\percent} vs.~\SI{33.5}{\percent}), a~zároveň výrazně nízkého pokrytí kolektivními smlouvami (\SI{32}{\percent}).
Dánský model --- nejnižší pracovní doba, nejnižší \ac{AROPE} bez extrémního Giniho --- je dosažen při vysoké odborové hustotě (\SI{67}{\percent}) a~\ac{KS} pokrytí (\SI{80}{\percent}): demonstruje, že flexibilita trhu práce a~sociální ochrana nejsou protiklady, pokud jsou vyváženy silnými institucemi kolektivního vyjednávání.
Rakousco kombinuje nejvyšší \ac{KS} pokrytí (\SI{98}{\percent}, erga omnes) s~nízkým \ac{AROPE} (\SI{17.7}{\percent}) a~nadprůměrnou mzdou: ukazuje, že i~bez skandinávské odborové hustoty lze dosáhnout solidní mzdové ochrany, pokud je rozšiřování \ac{KS} právně garantováno.
ČR a~Slovensko --- oba s~nízkým pokrytím --- pracují více hodin, vydělávají méně v~\ac{PPS} a~mají vyšší daňový klín na pracovníka: tento trojúhelník naznačuje, že česká konkurenceschopnost je stále postavena na modelu levné práce, nikoliv na kvalitní institucionální infrastruktuře trhu práce.

Srovnávací tabulka \ref{tab:flexicurity} shrnuje klíčové indikátory trhu práce pro šest zemí střední a~severní Evropy za rok 2023 a~odhaluje strukturální paradox české ekonomiky: ČR dosahuje vysoké zaměstnanosti (\SI{81.7}{\percent}), nízkého \ac{AROPE} (\SI{12.0}{\percent}) a~relativně komprimovaného Giniho koeficientu (24,4) --- avšak za cenu nejvyššího počtu odpracovaných hodin v~souboru (\SI{39.6}{\hour\per\week}), daňového klínu přesahujícího Dánsko (\SI{38.1}{\percent} vs.~\SI{33.5}{\percent}), a~zároveň výrazně nízkého pokrytí kolektivními smlouvami (\SI{32}{\percent}).
Dánský model --- nejnižší pracovní doba, nejnižší \ac{AROPE} bez extrémního Giniho --- je dosažen při vysoké odborové hustotě (\SI{67}{\percent}) a~\ac{KS} pokrytí (\SI{80}{\percent}): demonstruje, že flexibilita trhu práce a~sociální ochrana nejsou protiklady, pokud jsou vyváženy silnými institucemi kolektivního vyjednávání.
Rakousco kombinuje nejvyšší \ac{KS} pokrytí (\SI{98}{\percent}, erga omnes) s~nízkým \ac{AROPE} (\SI{17.7}{\percent}) a~nadprůměrnou mzdou: ukazuje, že i~bez skandinávské odborové hustoty lze dosáhnout solidní mzdové ochrany, pokud je rozšiřování \ac{KS} právně garantováno.
ČR a~Slovensko --- oba s~nízkým pokrytím --- pracují více hodin, vydělávají méně v~\ac{PPS} a~mají vyšší daňový klín na pracovníka: tento trojúhelník naznačuje, že česká konkurenceschopnost je stále postavena na modelu levné práce, nikoliv na kvalitní institucionální infrastruktuře trhu práce.

Srovnávací tabulka \ref{tab:flexicurity} shrnuje klíčové indikátory trhu práce pro šest zemí střední a~severní Evropy za rok 2023 a~odhaluje strukturální paradox české ekonomiky: ČR dosahuje vysoké zaměstnanosti (\SI{81.7}{\percent}), nízkého \ac{AROPE} (\SI{12.0}{\percent}) a~relativně komprimovaného Giniho koeficientu (24,4) --- avšak za cenu nejvyššího počtu odpracovaných hodin v~souboru (\SI{39.6}{\hour\per\week}), daňového klínu přesahujícího Dánsko (\SI{38.1}{\percent} vs.~\SI{33.5}{\percent}), a~zároveň výrazně nízkého pokrytí kolektivními smlouvami (\SI{32}{\percent}).
Dánský model --- nejnižší pracovní doba, nejnižší \ac{AROPE} bez extrémního Giniho --- je dosažen při vysoké odborové hustotě (\SI{67}{\percent}) a~\ac{KS} pokrytí (\SI{80}{\percent}): demonstruje, že flexibilita trhu práce a~sociální ochrana nejsou protiklady, pokud jsou vyváženy silnými institucemi kolektivního vyjednávání.
Rakousco kombinuje nejvyšší \ac{KS} pokrytí (\SI{98}{\percent}, erga omnes) s~nízkým \ac{AROPE} (\SI{17.7}{\percent}) a~nadprůměrnou mzdou: ukazuje, že i~bez skandinávské odborové hustoty lze dosáhnout solidní mzdové ochrany, pokud je rozšiřování \ac{KS} právně garantováno.
ČR a~Slovensko --- oba s~nízkým pokrytím --- pracují více hodin, vydělávají méně v~\ac{PPS} a~mají vyšší daňový klín na pracovníka: tento trojúhelník naznačuje, že česká konkurenceschopnost je stále postavena na modelu levné práce, nikoliv na kvalitní institucionální infrastruktuře trhu práce.

\input{texparts/python/flexicurity_table}



